% PAGES: 71-73 (printed pages 55-57; begins mid-page-55 with the "2.5"
% section heading, immediately after the end of the proof of Lemma 2.3
% transcribed in sec02_c3_1.tex)

\section{Linear solvers using the LU decomposition without pivoting}

Let $A$ be a nonsingular square matrix of size $n$, and let $b$ be a column
vector of size $n$. If the matrix $A$ has an LU decomposition without
pivoting $A=LU$, then, solving a linear system $Ax=b$ is equivalent to
solving
\[
LUx \;=\; b.
\]
This is the same as solving
\[
Ly \;=\; b
\]
for $y$, which can be done using forward substitution since the matrix $L$
is lower triangular, and then solving
\[
Ux \;=\; y
\]
for $x$, which can be done using backward substitution since the matrix $U$
is upper triangular.

The pseudocode for solving a linear system using the LU decomposition
without pivoting of a matrix can be found in Table~2.6.

\begin{table}[H]
\centering
\caption{Linear solver using LU decomposition without pivoting}
\fbox{\begin{minipage}{0.9\linewidth}
\textbf{Function Call:}\\
$x = \text{linear\_solve\_LU\_no\_pivoting}(A,b)$
\bigskip

\textbf{Input:}\\
$A$ = nonsingular square matrix of size $n$ with LU decomposition\\
$b$ = column vector of size $n$
\bigskip

\textbf{Output:}\\
$x$ = solution to $Ax=b$
\bigskip

\noindent
$[L,U] = \text{lu\_no\_pivoting}(A)$; \hfill\textit{// LU decomposition of $A$}\\
$y = \text{forward\_subst}(L,b)$; \hfill\textit{// solve $Ly=b$}\\
$x = \text{backward\_subst}(U,y)$; \hfill\textit{// solve $Ux=y$}
\end{minipage}}
\end{table}

The operation count for the pseudocode from Table~2.6 is as follows:
\begin{itemize}
\item[$\bullet$] $\dfrac{2}{3}n^3+O(n^2)$ for the LU decomposition of $A$; cf.\ (2.44);
\item[$\bullet$] $n^2+O(n)$ for the forward substitution for solving $Ly=b$; cf.\ (2.8);
\item[$\bullet$] $n^2+O(n)$ for the backward substitution for solving $Ux=y$; cf.\ (2.19),
\end{itemize}
for a total operation count of
\[
\left(\frac{2}{3}n^3+O(n^2)\right) + \left(n^2+O(n)\right) + \left(n^2+O(n)\right) \;=\; \frac{2}{3}n^3+O(n^2);
\]
see (10.86) in Section 10.2.3 for a proof of the last equality.

\subsection{LU linear solvers for tridiagonal matrices}

Solving linear systems corresponding to tridiagonal matrices having LU
decomposition is often needed for practical applications such as cubic
spline interpolation and the finite difference solution of the
Black--Scholes PDE. We discuss the solution to tridiagonal linear systems
here, and furthermore in section 6.3 when investigating efficient linear
solvers for tridiagonal symmetric positive definite matrices.

We begin by showing that the $L$ and $U$ factors from the LU decomposition
without pivoting of a tridiagonal matrix are bidiagonal matrices.\footnote{More
generally, the $L$ and $U$ factors from the LU decomposition without
pivoting of a banded matrix of band $m$ are also banded of band $m$; see an
exercise from the end of this chapter and Stefanica \cite{38} for details.}

Let $A$ be an $n\times n$ tridiagonal matrix, i.e., such that
\begin{equation}
A(j,k) \;=\; 0, \quad \forall\, 1\le j,k\le n \text{ with } |j-k|\ge 2,
\end{equation}
and assume that the matrix $A$ has an LU decomposition without
pivoting.\footnote{A strictly diagonally dominated tridiagonal matrix is an
example of a tridiagonal matrix with LU decomposition without pivoting; see
section 4.3.} Let $L$ and $U$ be the LU factors of $A$. We will show that
$L$ is a lower triangular bidiagonal matrix, i.e., the only entries of $L$
that can be nonzero are the main diagonal entries $L(i,i)$, for $i=1:n$, and
the lower diagonal entries $L(i,i-1)$, for $i=2:n$, and that $U$ is an upper
triangular bidiagonal matrix, i.e., the only entries of $U$ that can be
nonzero are the main diagonal entries $U(i,i)$, for $i=1:n$, and the upper
diagonal entries $U(i,i+1)$, for $i=1:(n-1)$.

Following step by step the LU decomposition without row pivoting from
section 2.4.1, we compute the entries from the first row of $U$ as
$U(1,k)=A(1,k)$, for all $k=1:n$, and the entries from the first column of
$L$ as $L(k,1)=A(k,1)/U(1,1)$, for $k=1:n$. Note that $A(1,k)=0$ if
$3\le k\le n$ and $A(k,1)=0$ if $3\le k\le n$, since $A$ is a tridiagonal
matrix; cf.\ (2.47). Then, $U(1,k)=0$ if $3\le k\le n$ and $L(k,1)=0$ if
$3\le k\le n$, and therefore the only possible nonzero entries from the
first row of $U$ and from the first column of $L$ correspond to an upper
triangular bidiagonal form for $U$ and a lower triangular bidiagonal form
for $L$.

Moreover, since
\[
U(1,2:n) = (U(1,2)\ \ 0\ \ \dots\ \ 0)
\quad\text{and}\quad
L(2:n,1) = \begin{pmatrix} L(2,1) \\ 0 \\ \vdots \\ 0 \end{pmatrix},
\]
it follows that
\begin{align*}
L(2:n,1)U(1,2:n)
&= \begin{pmatrix} L(2,1) \\ 0 \\ \vdots \\ 0 \end{pmatrix} (U(1,2)\ \ 0\ \ \dots\ \ 0) \\
&= \begin{pmatrix}
L(2,1)U(1,2) & 0 & \dots & 0 \\
0 & 0 & \dots & 0 \\
\vdots & \vdots & \dots & \vdots \\
0 & 0 & \dots & 0
\end{pmatrix}.
\end{align*}

Thus, updating the $(n-1)\times(n-1)$ lower right part of $A$ as in (2.36),
i.e.,
\[
A(2:n,2:n) \;=\; A(2:n,2:n) \;-\; L(2:n,1)U(1,2:n),
\]
only involves changing the value of $A(2,2)$ to
\[
A(2,2) \;=\; A(2,2) \;-\; L(2,1)U(1,2),
\]
which preserves the tridiagonal structure of the matrix $A(2:n,2:n)$.

The tridiagonal structure of the updated part of $A$ and the bidiagonal
structure of $U$ and $L$ will be further preserved as every row of $U$ and
every column of $L$ are computed.

For example, assuming that the updated form $A(i:n,i:n)$ of the matrix $A$
is tridiagonal after $i-1$ rows of $U$ and $i-1$ columns of $L$ are
computed, the $i$-th row of $U$ is computed as $U(i,k)=A(i,k)$, for all
$k=i:n$, and the $i$-th column of $L$ is computed as $L(k,i)=A(k,i)/U(i,i)$,
for all $k=i:n$. Since $A(i:n,i:n)$ is tridiagonal, it follows that
$A(i,k)=0$ if $i+2\le k\le n$, and therefore
\[
U(i,k)=0, \quad \forall\, i+2\le k\le n,
\]
and also that $A(k,i)=0$ if $i+2\le k\le n$, and therefore
\[
L(k,i)=0, \quad \forall\, i+2\le k\le n.
\]

Thus, the only possible nonzero entries from the $i$-th row of $U$ and from
the $i$-th column of $L$ correspond to an upper triangular bidiagonal form
for $U$ and a lower triangular bidiagonal form for $L$.
